\section*{Tree codes and palette-alternating models}
We turn to formally define tree codes. Let $T$ be a rooted binary tree that is endowed with an edge coloring from some ambient color set (or alphabet) $\Sigma$. Let $u,v$ be a pair of vertices in $T$ with equal depth and a least common ancestor $w$. Let $\ell$ be the distance, in edges, from $u$ to $w$. Let $p_u,p_v\in\Sigma^\ell$ be the sequences of colors on the path from $w$ to $u$ and to $v$, respectively. We define $h(u,v)$ to be the relative Hamming distance between $p_u$ and $p_v$.
\begin{authorclaim}{Definition1 (Tree codes [29])}
Let $T$ be the complete rooted infinite binary tree. The tree $T$, together with an edge-coloring of $T$ by a color set $\Sigma$ is called a tree code with distance $\delta$ if for every pair of vertices $u,v$ with equal depth it holds that $h(u,v)\ge\delta$. When there is no $\delta>0$ for which $T$ is a tree code with distance $\delta$ we say that $T$ has vanishing distance.
\end{authorclaim}
Schulman [29] proved that for every distance parameter $\delta<1$ tree codes with a constant number of colors $c=c(\delta)$ exist. Although tree codes are used in different ways by different interactive coding schemes, one aspect is common to all: When a party wishes to send a bit, a suitable color from $\Sigma$ is sent instead. Thus, the rate of all tree-code-based coding schemes is bounded above by $1/\log_2|\Sigma|$. One is led to ask a natural combinatorial question--what is the least number of colors in a tree code with non-vanishing distance?
\paragraph{1.1 Tree codes:4 colors suffice and are necessary}
We first observe that 3 colors do not suffice and, as a result, the rate of every tree-code-based coding scheme cannot exceed1/2, let alone approach capacity. Consider any 3-color tree code. First, we may assume that every two siblings are connected to their parent with edges having distinct colors as otherwise the distance of the tree code is 0. Let $u,v$ be any two vertices. Out of $u,v$ go four edges and so by the pigeonhole principal in every 3-coloring, two of these edges share the same color. By the above, one of these edges goes out of $u$ and the other goes out of $v$. This implies that, starting from the two sons of the root, one can construct two paths of any desired length $n\ge1$ with the same color pattern, establishing that the tree has vanishing distance.
Based on the ideas Schulman introduced to prove the existence of tree codes with a constant number of colors, we complement the above observation and establish that 4 colors suffice for a tree code with non-vanishing distance.
\begin{authorclaim}{Theorem2}
There exists a 4-color tree code with distance 0.136.
\end{authorclaim}
\paragraph{Definition3 (Palette-alternating tree codes)}
Let $\Sigma_0,\ldots,\Sigma_{c-1}$ be (not necessarily distinct) sets. Let $T$ be the complete rooted infinite binary tree. A palette-alternating tree code is an edge-coloring of $T$ where at layer $t\in\mathbb N$ the colors are taken from the set $\Sigma_{t\bmod c}$. $T$ is said to have distance $\delta$ if for every pair of vertices $u,v$ with equal depth it holds that $h(u,v)\ge\delta$. We define the rate $\rho$ of $T$ to be the number satisfying
\[
\frac1\rho=\frac1c\sum_{i=0}^{c-1}\log_2|\Sigma_i|.
\]
We suggest that the flexibility introduced by palette-alternating tree codes allows one to better capture the notion of rate in the online setting. Indeed, the importance of rate is only significant when "long" messages are being sent and so, informally, using a big palette of colors only once in a while should not be considered as an indication of poor rate. Our definition of rate formalizes that property. Note that we still insist on having the distance measured in terms of worst-case--a must as we wish to replace tree codes with palette-alternating tree codes in interactive coding schemes. It is only the rate that is being, in a sense, amortized.
As mentioned, we prove that palette-alternating tree codes can have rate approaching arbitrarily close to1 while maintaining non-vanishing distance, thus bypass the1/2 bound proven for (standard) tree codes.
\section*{2 Preliminaries}
Unless otherwise stated, all logarithms are taken to the base2. We denote by $\mathbb N$ the set of natural numbers (of course, including0), and write $\mathbb N_1$ for $\mathbb N\setminus\{0\}$. For integers $a\le b$ we write $[a,b]$ for all integers in this interval. For an integer $c\ge1$, we let $[c]=\{1,2,\ldots,c\}$. We follow the convention that strings are indexed starting from1. For two strings $x,y\in\Sigma_1\times\cdots\times\Sigma_n$, we denote by $\Delta(x,y)$ their hamming distance. We make use of the following standard inequalities.
\begin{authorclaim}{Lemma7}
For every integers $1\le k\le n$ with $k/n=\delta\le1/2$ it holds that $\sum_{i=0}^k\binom ni\le2^{H(\delta)n}$.
\end{authorclaim}
\begin{authorclaim}{Lemma8}
For every $0<x<1/2$ it holds that $x/(2\log_2(6/x))\le H^{-1}(x)\le x/\log_2(1/x)$.
\end{authorclaim}
